%!TEX root = ../dissertation.tex
\begin{savequote}[75mm]
[A formal program proof] is not merely a method to make absolutely sure we have not made a mistake in a proof, but also a tool that shows us and compels us to understand why a proof works.
\qauthor{Georges Gonthier}
\end{savequote}

\chapter{Implementation}\label{chapter implement}
Vizing's theorem relies heavily on a number of auxiliary data structures, including graphs, edge colorings, edge sets, fans, and Kempe chains. For each of these, we needed to provide a well-suited definition, state basic properties, and reason about how it interacts with other structures and functions under particular circumstances. Ideally, the properties should be readily extractable from the chosen representation, and we should minimize the complexity required to properly interact with the other structures. Thus, finding the right representations in our code was crucial to the success of this project. We divided our code into, roughly, one file per core data structure or type as well as a file called \fxn{vizing.v} which contains the final proof of Vizing's theorem, pulling from these other files. The list of files and related statistics can be viewed in Table~\ref{tab:stats}.
\begin{table}[htbp]
  \centering
    \begin{spacing}{1.0}
    \sffamily
    \begin{tabular}{lrrrr}
    \hline
    \textsc{File} & \textsc{Lines} & \textsc{Lemmas} & \textsc{Definitions} & \textsc{Admits} \\
    \hline\hline
    \texttt{alternate\_path} & 248 & 25 & 11 & (15) \\
    \texttt{aux} & 169 & 24 & 4 & 0 \\
    \texttt{edge\_coloring} & 486 & 34 & 17 & 1 \\
    \texttt{fans} & 330 & 37 & 12 & 2 \\
    \texttt{kempe} & 343 & 22 & 6 & 3 \\
    \texttt{spath} & 195 & 10 & 2 & 0 \\
    \texttt{vizing} & 214 & 4 & 0 & 0 \\
    \hline
    \textbf{Total} & \textbf{1985} & \textbf{156} & \textbf{52} & \textbf{6} \\
    \hline
  \end{tabular}
  \end{spacing}
  \caption{Per-file breakdown of the VizingThesis library. The lines of code are excluding blank/commented lines. The lemma count includes admitted proofs, lemmas, corollaries, propositions, and theorems. The proofs in \fxn{alternate\_path} are intentionally admitted, so we do not include them in our total.}
  \label{tab:stats}
\end{table}
Aside from \fxn{vizing.v}., each file listed centers around new definitions (e.g., a Kempe chain, fan, etc.) and functions to modify the structures (e.g., swap colors, rotate a fan). The majority of each file is devoted to lemmas about properties which are derived directly from the definition (e.g., all k-edge colorings are proper edge colorings), change with functional operations (e.g., the effect of swapping colors on an absent set), or describe interactions with objects in other files (e.g., performing a Kempe swap with a fan). The purpose of this chapter is to elaborate on some of the representations we chose, why these were chosen, and how these impacted the project. 

The remaining four sections are organized as follows: Section~\ref{section about rocq} provides background on the Rocq proof assistant before we dive into more technical details. Then Section~\ref{section library} shares how graphs are seen in Rocq and what supporting material is already available within Rocq's libraries. Section~\ref{section d&d} is the heart of this chapter; we split this into four subsections, each roughly focused on one of the code files from Table~\ref{tab:stats}. Finally, we conclude with Section~\ref{section admit it} by addressing the admitted proofs alluded to in Table~\ref{tab:stats}.

\section{Introduction to Rocq}\label{section about rocq}
Proof assistants and formal verification systems can exist because of the \textit{Curry-Howard correspondence}~\cite{CurryHow} which equates formal proof systems and type systems. Mathematical formulas and proofs can be interpreted as types and programs for models of computation, and vice versa. The formal system underlying the Rocq proof assistant utilizes the \textit{calculus of inductive constructions} (CIC), an extension of the Curry-Howard correspondence. With this mindset, Rocq can be understood as both a dependently typed functional programming language and as a formal verification system where formulas are types and proofs are programs. Just as every program has a type, every proof term is associated with a formula or proposition. We write $x : T\,$ to mean ``$x$ is of type $T.$'' This can also be read as ``$x$ belongs to $T$'' if we think of types as sets; in this sense, the term $x$ is said to \textit{inhabit} $T.$ A type is \textit{inhabited} if it contains at least one term. For most programming languages, we can interpret $x : T\,$ as ``$x$ is a program of type $T$.'' In Rocq, this could also be interpreted to mean ``$x$ is a proof term for the proposition $T$.'' Determining if some term $x$ is indeed a proof for a proposition $T$ is equivalent to deciding if $x : T\,$ is a well-typed expression. Rocq's kernel is a type checker, and it follows the set of typing rules outlined by the CIC (see \cite{cic_doc}) to verify programs and proofs.

As an introduction to how one might formalize a mathematical proof in Rocq, we will walk through an implementation of the following lemma. 
\begin{lemma}\label{example}
    Suppose $A$ and $B$ are two finite sets containing elements of some type $X.$ For any element $a$ of type $X,$ if we have $a \in A$ and $a \notin A \backslash B,$ then $a \notin B.$
\end{lemma}
This Lemma is named \fxn{notin\_setD} in Rocq and is the first proof in the \fxn{aux} file. It is meant to assist with arguments about sets in later files. In Rocq, we type \fxn{(X : finType)} to mean $X$ is any finite type, \fxn{(A, B : \{set X\})} to mean variables $A$ and $B$ are any sets containing elements of type $X,$ and \fxn{(a : X)} to mean $a$ is any element element of type $X.$ Set membership or exclusion is denoted with ``$\backslash $\fxn{in}'' and ``$\backslash$ \fxn{notin},'' and the characters ``:$\backslash$:'' refer to set difference. Then we can write the lemma as
\begin{minted}{coq}
Lemma notin_setD 
    {X : finType} (A B : {set X}) (a : X) (hA : a \in A) (hAB : a \notin A :\: B)
    : a \in B.
\end{minted}
where \fxn{hA} and \fxn{hAB} are arbitrary names we have assigned to our hypotheses. Thus, this lemma is of the form $x : T\,$ where $x$ is our list of variables and hypotheses and $T$ is the conclusion $a \in B.$ Just as with programs, we can write this in curried form:
\begin{minted}{coq}
Lemma notin_setD {X : finType} (A B : {set X}) (a : X) :
  a \in A -> a \notin A :\: B -> a \in B.
\end{minted}
This has the same meaning as above, though it now better resembles the corresponding mathematical statement: $a \in A$ implies $a \notin A \backslash B$ implies $a \in B.$ Note that this is logically equivalent to $a\in A$ and $a \notin A \backslash B$ implies $a \in B.$

To verify the lemma, we must enter a \textit{proof state} by stepping to this location in the code and construct the proof term to show our hypotheses have type $a \in B.$ A proof state consists of the context and the goal or goals. The \textit{context} displays local hypotheses or assumptions (e.g. $a : X\,$), and the \textit{goal} displays the current proof term we need to construct. Proof terms quickly become cumbersome and difficult to interpret. Like most proof assistants, Rocq supports \textit{tactics} to assist the programmer in building a proof term with standard logical reasoning techniques. Our code is written in the small-scale reflection language (SSReflect). This is a formal proof language based on Rocq and functions as an extension of Rocq's traditional tactics with the intent of shortening proofs and generalizing traditional tactics for a wider range of applications. Rocq's Mathematical Components library, graph theory library, and proof of the four color theorem are all written in SSReflect. An example of a tactic in SSReflect is the \fxn{rewrite} tactic. It supports rewriting two items given a proof that they are equivalent (e.g. rewriting an occurrence of $A$ with $B$ given $A = B$), similar to Rocq's traditional \fxn{rewrite} tactic. Additionally, it incorporates definition expansion and partial evaluation. Another distinction of SSReflect are the tacticals ``:'' and ``=>'' which can be attached to tactics and move hypotheses from the context to goal or vice versa, respectively, without requiring an additional \fxn{intros} tactic or similar. Every goal is encouraged to end with a terminating tactic, that is, a tactic that generates an error if it does not close the goal. In SSReflect, this is commonly done by preceding the last line of a proof with the ``\fxn{by}'' tactical. Not only does this convert any subsequent tactics into terminating tactics; it also solves any trivial goals remaining after the tactic is applied.

Given this background, we close this section with the proof of our example Lemma~\ref{example}. We present, side-by-side, the proof in Rocq and a traditional mathematical proof of the lemma. The exact library results we are rewriting with (\fxn{in\_setD}, \fxn{andbT}, and \fxn{negbK}) are not as important here as the general structure of a proof in Rocq, but they essentially align with the steps taken in the mathematical proof. 

\noindent
\begin{minipage}[t]{0.35\textwidth}
\textbf{Rocq proof}
\begin{minted}{coq}
rewrite in_setD.
move=> hA. rewrite hA.
by rewrite andbT negbK.
\end{minted}

\end{minipage}
\begin{minipage}[t]{0.55\textwidth}
\textbf{Mathematical Proof}
\medskip
\begin{spacing}{\dcompressedspacing}
By definition of set difference, the assumption $a \notin A \backslash B\,$ is equivalent to $\neg (a \notin B \wedge a \in A).$\\
\medskip
Since $a \in A,$ this is equivalent to $\neg(a \notin B \wedge \text{true}).$\\
\medskip
Therefore, we have $\neg(a \notin B)= a \in B.$ This is what we wished to show, so we are done.
\end{spacing}
\end{minipage}
\hfill
% Calculus of Constructions (CoC)\cite{coc}
% Unlike classical logic, the CIC does not assume the law of excluded middle. For any proposition $P,$ this law asserts either $P$ or its negation $\neg P$ must be true. Instead, the CIC and related theories follow a \textit{constructivist} logic, requiring the explicit construction of a witness to prove existence as opposed to a proof showing non-existence implies a contradiction. \note{implications for us}
\section{The Graph Theory Library}\label{section library}
Rocq's main library for formalized mathematics is the Mathematical Components library, or MathComp. MathComp is written in SSReflect, and it was used in Gonthier's proof of the four color theorem~\cite{4CRocq} as well as a proof of the odd order theorem~\cite{groupTheory}. While we did reference this library, especially for finite sets, sequences, and natural number arithmetic, the finite graphs in MathComp were built with planar graphs in mind and did not meet our needs. In 2019, Christian Doczkal and Damien Pous released a more general graph library that can now be installed as a Rocq package~\cite{Doczkal, graphLibRock}, and we build our project off this. The library already includes Hall's Marriage theorem, Menger's theorem, and results related to vertex colorings, matchings, treewidth, connectivity, and a few other areas. They begin with a directed graph, called an \fxn{digraph}, as a structure consisting of a finite type $G$ and a relation $R: G \to G \to$ \fxn{bool} where \fxn{bool} is an inductive data type that can either be true or false. The vertices of the graph are the inhabitants of $G$ and the relation $R$ denotes edges. For two vertices $x, y : G,$ if the relation on $x$ and $y$ is true, then there is an edge between $x$ and $y$ in the graph, and this is denoted as $x$ \fxn{--} $y.$ A simple graph, called an \fxn{sgraph}, is a \fxn{digraph} with the a symmetric and irreflexive restriction imposed on the edges. That is, a simple graph is a directed graph in which we have an edge $x$ \fxn{--} $y$ if and only if we have an edge $y$ \fxn{--} $x$ and there can be no self-loops ($x \neq y$). Sometimes we will refer to edges with this notation, and sometimes we will refer to them as sets with elements of type $G,$ e.g. $[x; y].$ The latter does not depend on the relation $x$ \fxn{--} $y$ to be true, for we can easily write $[x; y]$ for any two vertices in the graph, even if there is no edge between $x$ and $y.$ However, sometimes the existence of an edge is unnecessary or it is implied by other hypotheses; in these cases, it is sometimes easier to work with the set notation rather than the relation.

\section{Implementation details and decisions}\label{section d&d}
\subsection{Edge Colorings}
The first hurdle in our implementation was defining edge colorings. This foundational step was nontrivial and deserved sufficient deliberation. The main contender was rooted in set theory. In the library, vertex colorings are seen as partitions of the vertices into disjoint sets. Matching this form, we considered `coloring' the edges by partitioning them into disjoint `color' sets themselves. While this promotes the relationship between edge colorings and matchings, we ultimately rejected it for a number of reasons. Firstly, we would likely want to adjust the definition’s structure because the default approach introduced a cumbersome nested set notation. Additionally, it does not as cleanly lead to an extraction of the `color' assigned to a given edge nor to reasoning about non-proper edge colorings. Thus, we rejected this viable alternative, though it could be interesting to compare our product to such an implementation. 

Before migrating to a functional approach, we briefly considered constructing a new \fxn{CGraph} type which, similar to the library’s coercion between simple and directed graphs, would extend the \fxn{SGraph} with a symmetric coloring relation over the edges. Another option would be to modify a vertex coloring to assign \emph{two} distinct colors to each vertex and require at least one color be shared between any two vertices with an edge relation. While these and many others were enjoyable and, in some cases, enlightening to ponder, most proved too inconvenient or bulky to be practically worthwhile.

Ultimately, we converged on the simple mapping 
$$ C_{G, Colors} : E(G) \to Colors$$
which sends an edge in some graph $G$ to a color in some set $Colors$ of valid colors. In Rocq, this dependent mapping is defined as follows:
\begin{minted}{coq}
Definition edgeColoringType (G : sgraph) (ColorType : finType) : Type := 
    {set G} -> ColorType.
\end{minted}
Thus, \fxn{edgeColoringType} is the type for a map sending an edge (represented as set of vertices) to a `color' where the color is some object of type \fxn{ColorType}, and it is dependent on a simple graph \fxn{G} and finite type \fxn{ColorType}. We intentionally granted flexibility in the \fxn{ColorType} rather than avoid the additional argument but limit ourselves to something less general such as natural numbers. In fact, we almost never need to actually specify a finite type to be the \fxn{ColorType}. We only explicitly state our \fxn{ColorType} when we create an injective coloring to prove Vizing's lower bound, matching each edge to a unique color. In this case, the \fxn{ColorType} is the type of the edges themselves (\fxn{\{set G\}}) which is valid because $G$ itself is a finite graph.

Once we had established edge colorings, we added methods to recolor edges, swap edge colors, and extend an edge coloring to incorporate new edges. Swaps can be thought of as recoloring twice, so they were relatively similar to the recoloring implementation. We checked our definition of recoloring by proving, in one line each, that the recolored edge maps to its new color and for all other edges, the color of that edge remains the same. In other words, if we let $D = C[e \mapsto c_0]$ be the edge coloring with some edge $e$ recolored as some color $c_0$, then we show $D(e) = c_0$ and $D(f)  = C(f)$ for all edges $f \neq e$ as expected. Recoloring itself receives little direct attention in mathematical literature, in part because it is straightforward in nature and its applications to other more complicated structures or operations are of greater interest. However, a couple of recoloring propositions fundamental to Vizing's theorem required some level of casework. This casework is often hidden within a succinct mathematical justification, but proof assistants ensure no case is accidentally missed. The last statement of Misra and Grieg's algorithm is
\begin{quote}
``Rotation leaves [the color] $d$ free at [the vertex] $w$, since the only edge incident on $w$ that has its color changed by the rotation is [the edge] $(x, w)$ and this edge is uncolored by
the rotation. Since $d$ is also free at $x$, and since the rotation uncolors $(x, w)$, giving $(x, w)$ the color $d$ increases the number of colored edges and leaves the graph valid,''~\cite{Misra}.
\end{quote}
The lemma \fxn{recolor\_proper} corresponding to the last claim alone (the recoloring ``leaves the graph valid'') accounts for 30 lines of code excluding the assumptions and proof statement itself, for the Rocq kernel is well-suited to handle casework but must be told which cases to consider. It cannot trivialize a logical assertion at the same level of abstraction as Misra and Grieg can expect from their readers. As was the case throughout the implementation process, we had to intimately understand \textit{why} each statement was true in order to validate it in Rocq.

To extend an edge coloring $C$ when an edge was added to a graph $G$, we built a \textit{de facto} partial coloring $C\,'$ from $C:$
$$C\,'(e) = \begin{cases}
Some (C(e)) & e \in E(G)\\
None & e \notin E(G)
\end{cases}$$
Instead of creating an uncolored edge, we assigned the color `None' to the new edge. All other colors were wrapped inside a `Some' to ensure there were no conflicts (an edge in $E(G)$ may have already had color `None'). We then proved extending an edge coloring preserved the properness of a coloring and added exactly one new color to the color set. Although these seem like natural conclusions of the definition, the edge coloring $C$ and the extended edge coloring $C\,'$ have different types that became challenging to manage, especially because we already were dealing the graph $G$ and the graph $G$ with an edge added, which also have different types. In particular, we struggled to build fans starting at the newly colored edge for both extended edge colorings and extended k-edge colorings. Eventually, we decided to have two separate methods of constructing fans. Vizing's theorem only relied on the fan constructed from an extended k-edge coloring, so this was sufficient. However, this may indicate we may be missing a feature of Rocq or there is limitation in our representation, and investigating this could be a direction for future work.

Proper edge colorings, referred to as \fxn{properEdgeColoringType}, and k-edge colorings, referred to as \fxn{kEdgeColoringType}, followed directly from the \fxn{edgeColoringType}. As noted in Chapter~\ref{chapter math}, these are technically the same by our mathematical definitions. In the code, we utilize the \fxn{kEdgeColoringType} when we wish to hide the \fxn{ColorType} and emphasize the cardinality of the color set. This occurs when we define the chromatic index because it is not be dependent on the choice of \fxn{ColorType} and `chi’ or `k’ is more concise than, for instance, typing \fxn{\#|c[E(G)]|} to mean the cardinality of the image of an edge coloring $c.$ As such, unlike the \fxn{edgeColoringType} and \fxn{properEdgeColoringType}, the \fxn{kEdgeColoringType} is dependent on a graph \fxn{G} and a natural number \fxn{k} but not some finite type \fxn{ColorType}. We included several results regarding lower and upper bounds of \fxn{\#|c[E(G)]|} and created an injective coloring for any graph to assign every edge exactly one color, though we omit discussion of these here for brevity. 

We then say a graph $G$ is k-edge colorable if the type \fxn{(kEdgeColoringType G k)} is inhabited given some graph $G$ and natural number $k$. Recall from Section~\ref{section about rocq} that a type is inhabited if it contains at least one term. Determining if a type is inhabited is undecidable in general, and to convince Rocq a type is inhabited we must explicitly construct a term of that type. In this case, a graph $G$ is k-edge colorable if we can show we have a term of type \fxn{(kEdgeColoringType G k)}. A natural number $\ci$ is the chromatic index of a graph if the graph is $\ci$-edge colorable and for all $k$ less than $\ci$ the graph is not $k$-edge colorable. Hence, once we move to proving the upper bound of Vizing's Theorem, our goal is to construct a term of type \fxn{(kEdgeColoringType G ($\Delta$(G) + 1))}. To do this, we required the assistance of Kempe chains, terminating functions to construct maximal-length structures, and a few new propositions about the edges of graphs.

\subsection{Kempe Chains}\label{section Kempe implement}
For the Kempe chain, the literature suggested two main representations: a recursively extended alternating path and a connected subgraph containing a given vertex. We focus on the latter, though we also explore the former as a means of comparison and deeper understanding. The former begins with a length-one path, adding edges of alternating color until it is no longer possible. Following Doczkal et. al.’s packaged paths detailed in \cite{Doczkal}, we could quickly prove statements about extending and merging paths. Similar to fans, the construction of a maximal alternating path required a proof of termination. However, termination here relies on the properties of a proper edge coloring to avoid returning to a vertex already in the path as we extend it. We can push this more complicated proof away from the recursive function itself (and allow the function to be applicable even for improper edge colorings) by requiring the vertices to be unique. That is, the function terminates if no edges adjacent to the last vertex have the correct color or if such an edge exists but is adjacent to a vertex already in the path. The proof of termination also now nearly matched the maximal fan case, discussed in the next section. We would then later need to prove that the second case is unreachable when the relevant edge coloring is proper. The theorem \fxn{Vizings\_altpath} in the file \fxn{vizing.v} proves Vizing's theorem with Kempe chains implemented as these alternating paths, and they are defined in the file \fxn{alternating\_path.v}. Most of the supporting lemmas as a proof-of-concept are stated but admitted. Bhoja follows this representation in Lean and the curious reader can reference his paper \cite{Bhoja} and code \cite{bhojaCode} for completed proofs. We also added an edge membership predicate for paths to support the alternating paths. This is fully implemented in \fxn{spath.v}, and Section~\ref{section additional} discusses it in greater depth. 

Meanwhile, our final product supports Kempe chains interpreted as connected components of tailored subgraphs. Two main reasons ultimately led us to pursue this representation. Firstly, subgraphs and graph components already existed in the library. Modeled especially after the induced and complement graph types, we began by characterizing the edges of our Kempe subgraph with a graph $G,$ an edge coloring $C,$ and two colors $c_1$ and $c_2.$ If we let $H$ be the Kempe subgraph, then the edges of $H$ were defined as a subset of the edges of $G:$ 
$$
E(H) = \{e \mid e \in E(G) \wedge (C(e) = c_1 \vee C(e) = c_2)\}
$$
This translated to the following in Rocq:
\begin{minted}{coq}
Definition Kempe_rel := 
    [rel x y : G | x -- y && ((C [set x; y] == c1) || (C [set x; y] == c2))].
\end{minted}
That is, there is an edge relation between any two vertices $x$ and $y$ if there is an edge relation between $x$ and $y$ in $G$ and the edge $(x, y)$ is colored either $c_1$ or $c_2$ under $C.$ We could then define the subgraph $H$ of $G$ with the vertices $V(H) = V(G)$ and edges $E(H)$ from the symmetric, irreflexive edge relation \fxn{Kempe\_rel}. We were required to prove basic properties of the edge relation and resulting subgraph in order to properly define the subgraph as a simple graph and be able to manipulate it later on. Yet, we could closely mirror the proofs for other subgraphs in the library, and, though it demanded some time, was not particularly challenging. We also verified that our implementation acted as expected. The most direct lemma in this process, labeled \fxn{mem\_kempe}, stated $$e \in E(H) \iff e \in E(G) \wedge (C(e) = c_1 \vee C(e) = c_2)$$ and turned out to be quite useful -- we used it with the rewrite tactic 14 times. 

Graph components were already in the graph library, so a Kempe chain was defined as the component of the Kempe subgraph $H$ containing some specified vertex $v.$ This leads us to the second benefit of the subgraph representation: we avoided the tricky termination proof associated with the maximal Kempe chain. A component of a graph is immediately maximal; it by definition contains all edges connected to the specified vertex. Thus, we could reason directly with the component and degrees of its vertices. We still had to show that a Kempe chain was an acyclic path in order to restrain the number of distinct single-degree vertices. Originally, we expected our proof to apply the definitions of paths and cycles to disprove the existence of disconnected components and branching. Eventually though, we settled on Proposition~\ref{chain_two_endpts} and related results due to complications and limitations discussed in Section~\ref{section Kempe definition} of Chapter~\ref{chapter math}.

The Kempe chains themselves were the more challenging half of this section. Kempe swaps were akin to the recolor and swap operations, so much of the high-level logic could be applied. Some proofs were slightly longer because we also had to reason about edges in or not in the chain. The lemmas \fxn{not\_kempe\_edge} and \fxn{is\_kempe\_edge} assisted here and served to check assumptions about how the inversion should behave. They stated the color of an edge in the inverted edge coloring given its initial color and whether it was in the chain. Most of the other results on Kempe swaps are stated as lemmas in Section~\ref{section Kempe definition}.

\subsection{Recursion and Termination}
All functions in Rocq must terminate. By default, the kernel checks a \textit{guard condition}, ensuring recursive calls are made on structurally smaller arguments~\cite{wf_doc}. The author can explicitly label the decreasing argument when necessary. However, sometimes this is not enough, as was the case with maximal fans and maximal alternating paths. The claim that, given some object, there exists a maximal form of this object can be easily overlooked, taken as a fact rather than an assumption in handwritten proofs. The beauty of proof assistants is that no assumption can be unintentionally overlooked.

\fxn{Program Fixpoints} and \fxn{Equations} exist to assist in proving more complicated cases of dependent pattern matching and structural or well-founded recursion. As we discovered when attempting to practice with the \fxn{Program} mechanism, \fxn{Program} relies on an additional axiom -- JMeq, or John Major's Equality. The details on this axiom are out of the scope of this thesis, but we generally leaned against assuming additional axioms if not necessary. Luckily, the \fxn{Equations} plugin has no additional axioms, generates a functional elimination principle for inductive reasoning, and has, in our opinion, cleaner syntax. 

To actually define terminating recursive functions for the fans and Kempe chains, we found a lower-bounded value that decreased at each recursive call. The vertices within a fan must be neighbors of the central vertex and must be unique, so each time a new vertex is added to the fan, the number of neighbors of the central vertex that are not already in the fan decreases. The vertices of a Kempe chain must also be unique, so each time a new vertex pair is added, the number of edges in the graph that are not already crossed by the chain decreases. Note that these values were bounded and finite because we only considered finite graphs. 

These functions also required constructors to build new fans and paths. Each additional vertex created a longer sequence, but we needed to create a longer fan or a longer path. Fans and paths are sequences with additional constraints, so we wrote tailored sequence cons functions that would return not just a sequence but a fan or path, respectively, if the corresponding requirements were met. This allowed the recursive maximal fan and path functions to return objects with the expected type. Finally, referring to Bhoja~\cite{Bhoja}, we formalized the notions of maximality for our fans and paths and proved that the fans and paths returned by the recursive functions were indeed maximal.
\subsection{Edge Operations}\label{section additional}
As we converted our written proofs into code, we sometimes discovered basic concepts (especially related to edges) such as the edge set of a vertex had not yet been integrated into the library. Thus, we created files \fxn{aux.v} and \fxn{spath.v} to contain broadly applicable material that would be a valuable addition to the graph theory library independent of Vizing's theorem. This material includes, but is not limited to, the edge set of a vertex, edge membership in paths, additional results about the interior of a path, the maximum degree of a graph, and a special case of the lemma \fxn{bigmax\_leq\_pointwise} in the preliminaries of the graph library when exact equality is known. For the sake of brevity, we will only discuss the most interesting of these here: edge sets and path membership. 

The (open) neighborhood $N_G(x)$ of a vertex $x$ in some graph $G$ is the set of vertices $y$ in $G$ such that vertices $x$ and $y$ are adjacent in $G.$ On the other hand, the edge set $E_G(x)$ of a vertex $y$ in some graph $G$ is the set of edges $\e{x}{y}$ in $G$ such that vertex $x$ is adjacent to vertex $y$ in $G.$ Using set notation, we can express these formally:
\begin{align*}
    N_G(x) &= \{y \mid (x, y) \in E(G) \}\\
    E_G(x) &= \{(x, y) \mid y \in N_G(x)\}
\end{align*}
The graph library already had established neighborhoods, so we only needed to define edge sets of a vertex. We present both definitions here, adjusting variable names for comparison:
\begin{minted}{coq}
open_neigh (G : diGraph) (x : G) := [set y | x -- y].
edge_neigh (G : sgraph) (x : G) := [set [set x; y] | y in (open_neigh G x)].
\end{minted}
Both of these definitions use Rocq's set notation and can be read analogously to the mathematical definitions we provided above. One may notice that $G$ is type \fxn{diGraph} (a directed graph) for the open neighborhood and type \fxn{sgraph} (a simple graph) for the edge set. The scope of our project did not include directed graphs, so we did not concern ourselves with functionality beyond simple graphs, though, notably, simple graphs can be considered a special case of directed graphs. In certain scenarios, such as this definition, the two are essentially interchangeable, though not in all, such as when we rely on the symmetry of a simple graph. In any case, expanding our code to directed graphs where applicable should not be incredibly complicated and is a useful direction for future work. The main distinction we would like to point out is the use of the relation ``\fxn{x -- y}'' versus the set ``\fxn{[set x; y]}.'' In this library, the notation \fxn{x -- y} denotes the edge relation between two vertices $x$ and $y.$ The set notation \fxn{[set x; y]} ignores the relation, only storing the two vertices themselves. Yet, no information is lost because we know that a pair of vertices $x$ and $y$ are in $E_G(x)$ only if the vertex $y$ is in $N_G(x)$ and from this we can extract the relation. In fact, we prove 
\begin{minted}{coq}
 Lemma edgesSetP {G : sgraph} (x : G) (e : {set G}) :
    reflect (exists y, e = [set x; y] /\ x -- y) (e \in (edge_neigh G x)).
\end{minted}
to explicitly use this fact when necessary. For the majority of cases, the set notation was easier to reason about, especially given its alignment with the library's definition for $E(G),$ the set of all edges in the graph. This definition and a number of helpful lemmas can be found in the \fxn{aux} file of \cite{my_code}.

We also heavily reference definition, lemmas, and the general structure of the packaged paths designed for the graph theory library. We observed no consistent definition for edge membership of these paths. Kempe chains as alternating paths necessitate a definition of edge membership, so we took this as an opportunity to contribute a potential implementation. To determine if a path crossed some edge, we needed to check if the edge's two endpoint vertices appeared consecutively in the path. We found the best way to go about this was to mimic the definition of the path itself from MathComp. MathComp paths consist of a sequence of elements and a boolean predicate asserting some progression of the sequence. It also includes the view \fxn{pathP} directly relating a path to a proposition. We immediately derived our edge predicate from \fxn{pathP} by replacing the ``forall'' with an ``exists'' and dropping the boolean relation. Thus, given a path $P$ of a sequence of vertices $x_1, \dots, x_n$ and an edge $(u, v),$ we say the edge $(u, v)$ is in $P$ if the following logical statement is true:
$$
\exists i < n.\, (p_i = u \wedge p_{i+1} = v)
$$
The head of a path is explicitly stated in Rocq, so the path $P$ above would be written as $x_1 :: p$ where $p$ would be the sequence $x_2, \dots, x_n$ (which may be empty) and ``::'' is syntactic sugar for a cons operation, appending $x_1$ to the beginning of the sequence $p$. To access the i$^{\text{th}}$ index in a path, we use the \fxn{nth} function. This takes in a default value in case the i$^{\text{th}}$ index does not exist (e.g. the value of $i$ is a greater the length of the path itself). With this context, we present our raw edge membership predicate here.
\begin{minted}{coq}
Definition pathp_edge {G : sgraph} (x1 u v : G) (p : seq G) : Prop :=
    exists i, (i < size p) && (nth x1 (x1::p) i == u) && (nth x1 p i == v).
\end{minted}
The file \fxn{spath} also includes basic facts and reasoning about common path operations in relation to path edges. For example, the lemma \fxn{is\_path\_edge} asserts that if edge $(u, v)$ is a path edge, then the edge $u$ \fxn{--} $v$ is in the graph and the lemma \fxn{edgep\_path\_edge} proves that given a two-vertex path $P = x_1, x_2,$ the edge $(u, v)$ is an edge in $P$ if and only if $u = x_1$ and $v = x_2,$ or symmetrically, $v = x_1$ and $u = x_2.$ The longest of these is \fxn{cat\_path\_edge} which allows us to split a path, checking edge membership on the two sub-paths instead. These lemmas solidify the chosen representation and establish a solid foundation for future applications.
\section{Omitted Material}\label{section admit it} 
While we have proved over 100 lemmas and theorems in our code, there are still six proofs which end with \fxn{Admitted.} Like the a \fxn{sorry} in Lean, the \fxn{admitted} command terminates a proof even if there are open goals. In this section, we discuss why each of these were left incomplete and what would be required to finish them.
\begin{itemize}
    \item \textit{\fxn{extended\_absent}}: Suppose $C$ is an edge coloring of a graph $G,$ and suppose we add a new edge to $G$ to get the graph $G\,'.$ Let $C\,'$ be the extended edge coloring $C$ for $G\,'.$ This lemma states that for any vertex $x \in V(G) = V(G\,'),$ if there exists some color $c_0$ in $A_C(x)$ then the color $Some(c_0)$ is in $A_{C\,'}(x).$ We were somewhat unsatisfied with the way in an extended edge coloring interacted with our other definitions, so we did not yet complete this proof. It likely can be finished by expanding the definitions of absent sets, applying set membership and set difference predicates, and performing casework. Since we were limited in time, confident on the correctness of the statement, and it is not a crucial or striking result, we leave it as future work.
    \item \textit{\fxn{imset\_rot\_del\_edge}}: Suppose $C$ is an edge coloring of a graph $G$ and $R$ is the edge coloring $C$ after the rotation of some fan in the graph. Then this lemma can be seen as an extension of Lemma~\ref{imset_rot} which shows $C[E(G)] = R[E(G)]$. If $e_1$ and $e_2$ are any two edges in some graph $G$ such $C(e_1) = R(e_2)$ (where $e_1$ and $e_2$ may be the same edge), it states that $C[E(G) - e_1] = R[E(G) - e_2].$ That is, if the colors are the same in the respective colorings, then we can ignore the edges when considering the color set of the respective colorings and these will still be equivalent color sets. The proof, like most of our proofs with the rotation algorithm, utilizes induction. We have proven the two base cases, which includes proving the corresponding lemma for swapping edge colors, and have made progress on the inductive step. At this point, the proof could become messy as we reason about set membership, rotation, swapping edge colors, and the inductive hypothesis together. There may be a better way to approach this proof or phrase the lemma, and we leave this as future work.
    \item \textit{\fxn{rot\_proper}}: This corresponds to Lemma~\ref{rot_proper} from Chapter~\ref{chapter math}. That is, it shows the properness of an edge coloring is preserved when a fan is rotated. Again, this proof uses an induction argument. Now, we would need to incorporate a number of auxiliary lemmas such as \fxn{rot\_w0\_prop}, \fxn{rot\_notin}, \fxn{absent\_edge}, and \fxn{swap\_proper\_vertex}. We were unable to finish the proof in Rocq, but these auxiliary lemmas are implemented, and the remainder of the proof should closely follow the written proof in Chapter~\ref{chapter math}. 
    \item \textit{\fxn{in\_inverted\_absent}} and \textit{\fxn{notin\_inverted\_absent}}: These two lemmas are statements regarding the absent set of vertices that are in and not in a given Kempe chain, respectively, before and after the inversion occurs with a Kempe swap. These are intuitive statements that follow almost immediately from the definitions of Kempe chains and Kempe swaps. However, the proofs themselves would be on the longer side due to the casework required. Since we were limited on time, confident on the correctness of these statements, and these were not critical results, we left them as future work.
    \item \textit{\fxn{inverted\_fan}}: This proposition is meant to capture Proposition~\ref{inverted_fan} from Chapter~\ref{chapter math} describing a fan before and after we perform a Kempe swap, inverting the colors in a Kempe chain. Applying this proposition is a major step, and most of it is proven in the code; there are already 70 lines of code devoted to proof itself (excluding the statement and assumptions). It is incomplete because we have not yet proven that property~\ref{abs_prop} of fans holds for the sub-fan under consideration. To do so, we would need to better familiarize ourselves with MathComp's path type. This includes learning to separate paths into connected sub-paths and splitting the path property into these sub-paths. Additionally, we likely will need the previous two lemmas \fxn{in\_inverted\_absent} and \fxn{notin\_inverted\_absent}. Thus, proving this lemma should be a manageable goal but will take a little bit more time than some of the other admitted statements.
\end{itemize}